\section{Introduction}

Longitudinal imaging is the cornerstone of oncological response assessment. With cancer incidence projected to rise 47\% by 2040~\cite{bray2024global} and CT examination volumes increasing steadily~\cite{rocholl2025unstable}, radiologists face growing pressure to evaluate serial scans efficiently. Treatment response evaluation, typically governed by RECIST 1.1 guidelines~\cite{recist}, requires solving two distinct subtasks per lesion: \textit{retrieval}, identifying the same structure in a follow-up scan, and \textit{delineation}, measuring its volume change. Both remain predominantly manual, causing substantial reading time and inter-observer variability~\cite{hering2024improving,jacobs2021assisted,lu2021randomized}.

\noindent Existing automated approaches address this problem in complementary yet incomplete ways. \textbf{Point-only trackers}~\cite{yan2022sam,vizitiu2023multi} retrieve lesion centers, however, without volumetric delineation. \textbf{End-to-end trackers}~\cite{Rokuss_2025_CVPR} automate both retrieval and segmentation using a single baseline click, but operate as "black boxes": if the model tracks an incorrect structure or misses a splitting lesion, the resulting segmentation is clinically invalid with no opportunity for correction. \textbf{Decoupled registration-segmentaion pipelines}~\cite{hering2021whole} would permit verification of the propagated point but suffer twofold: (1) registration errors frequently displace prompts beyond the tolerance of segmentation models trained on well-centered clicks~\cite{rocholl2025unstable,scribbleprompt,nnInteractive,du2024segvol}, and (2) treating follow-up scans in isolation discards the \textit{longitudinal prior} (baseline appearance) needed to resolve ambiguous findings. Finally, \textbf{automated longitudinal models}~\cite{longiseg, szeskin2023liver, wu2023coactseg} leverage this prior lesion appearance for ehnanced segmentation but lack promptability, preventing user interaction or correction at all.\\


\noindent To bridge these methodological gaps, we argue that clinical deployment requires satisfying three criteria simultaneously: (i) explicit use of the baseline lesion as a longitudinal prior, (ii) a clinician-correctable mechanism to guarantee correspondence when tracking fails, and (iii) robustness to off-center prompts caused by registration or rapid correction. Because existing methods satisfy at most two, we propose a novel \textit{Verified Tracking} paradigm: registration proposes a follow-up location, a clinician verifies/corrects it, and the model segments using both the verified prompt and baseline context. This eliminates retrieval failures through minimal oversight, freeing the model to focus entirely on longitudinally-informed delineation. We present a unified framework for this workflow, combining early pixel-level prompt fusion~\cite{nnInteractive} with latent-space temporal difference weighting~\cite{longiseg}. Critically, we identify \textit{synthetic longitudinal pretraining} as a decisive enabler: without sufficient multi-timepoint data, longitudinal architectures collapse to single-timepoint shortcuts, ignoring the baseline scan entirely. Finally, to address longitudinal data scarcity, we curate and publicly release \textit{PanTrack} as a dedicated out-of-distribution benchmark. As public datasets with consistent, lesion-level instance annotations across multiple timepoints remain largely unavailable, this release fills a major gap, providing a valuable new resource for lesion tracking model development. Our key contributions are:

\begin{enumerate}
    \item \textbf{Verified Tracking formulation}: We formalize a workflow where follow-up prompts are registration-proposed and optionally corrected, offering a clinically safe middle ground between automation and control.
    \item \textbf{Longitudinal promptable segmentation model}: A unified architecture combining early prompt fusion and latent temporal difference weighting. Enhanced by promptable large-scale synthetic pretraining, our method won the MICCAI autoPET IV challenge, outperforming the state-of-the-art in automatic and verified tracking.
    \item \textbf{The \textit{PanTrack} benchmark}: To address multi-timepoint data scarcity, we publicly release 161 curated longitudinal CT scans (45 pancreatic cancer patients) to provide a rigorous out-of-distribution testbed for cross-domain tracking generalization and a novel model development resource.
\end{enumerate}

